\clearpage
\VerdierPhysicalPageAnchor{0188}%
\begingroup\small\setlength{\parskip}{0pt}\setlength{\emergencystretch}{1.5em}
\setlength{\abovedisplayskip}{4pt plus 2pt minus 2pt}%
\setlength{\belowdisplayskip}{4pt plus 2pt minus 2pt}%

\VerdierSemanticAnchor{ch3d-2-1-3-properties-cont}
一个初等态射。初等态射在基变换下稳定。设

\VerdierSemanticAnchor{ch3d-2-1-3-bi-infinite-chain}
\[
\cdots\longrightarrow X_n\xrightarrow{u_n}X_{n-1}
\xrightarrow{u_{n-1}}X_{n-2}\longrightarrow\cdots
\]
是一个在两个方向上都无界的复形态射列，并且对每个整数 $n$，态射 $u_n$
都是高度为 $-n$ 的初等态射。此列在
$\operatorname{comp}(\mathcal A)$ 中有射影极限与归纳极限（它在每个给定次数的
分量上所诱导的列在两个方向上都是定常的）。

\VerdierSemanticAnchor{ch3d-2-1-4}
\medskip\noindent\textbf{引理 2.1.4.}\quad
\textit{设 $m\colon X\longrightarrow Y$ 是
$\operatorname{comp}(\mathcal A)$ 中的一个拟同构。存在态射列}
\[
u_n\colon X_n\longrightarrow X_{n-1},\qquad n\in\mathbb Z,
\]
\textit{使得对每个整数 $n$，$u_n$ 都是高度为 $-n$ 的初等态射；并且存在态射}

\VerdierSemanticAnchor{ch3d-2-1-4-limits}
\[
\varinjlim_n X_n\xleftarrow{\ \sim\ }Y,
\qquad
\varprojlim_n X_n\xleftarrow{\ v\ }X,
\quad v\text{ 为同伦等价（第一章，2.5.11），}
\]
\textit{使得图}

\VerdierSemanticAnchor{ch3d-2-1-4-square}
\begin{center}
\begin{tikzcd}[column sep=5.6em,row sep=3em]
X \arrow[r,"v"] \arrow[d,"m"']
  & \displaystyle\varprojlim_n X_n \arrow[d] \\
Y \arrow[r,"{\sim}"'] & \displaystyle\varinjlim_n X_n
\end{tikzcd}
\end{center}
\textit{在同伦意义下交换。}

\VerdierSemanticAnchor{ch3d-2-1-4-proof}
首先注意，对每个态射 $m\colon X\longrightarrow Y$，若将其视为
$\operatorname{comp}(\mathcal A)$ 中的态射，则都存在一个满态射
$m'\colon X'\longrightarrow Y$ 以及一个在同伦意义下交换的图：

\VerdierSemanticAnchor{ch3d-2-1-4-epimorphism-diagram}
\begin{center}
\begin{tikzcd}[column sep=3.5em,row sep=3.2em]
X \arrow[rr,"v"] \arrow[dr,"m"'] && X' \arrow[dl,"{m'}"] \\
& Y &
\end{tikzcd}
\end{center}
\endgroup
